\mychapter{3}{Techno-economic modeling of NETs}
In order to develop optimal portfolios of negative emission technologies, the four relevant approaches were systematically modeled from a technological and economic perspective. The goal of these techno-economic models was to calculate the country-level availability and cost of these NET's in order to construct a global marginal abatement cost curve (MACC). These MACC's serve as the answer to the first research question, and, more importantly, provide the basis for the optimization models used to generate optimal portfolios of negative emission technologies delivering the required carbon removals.

\section{Methodology}
Generally, a bottom-up process was used to construct techno-economic models of the four relevant NET's. For each NET, the entire process from feedstock input to the final output of the carbon credits was modeled. This included the capital investments required to construct exemplary projects or plants, as well as the fixed and variable costs of operating these plants, primarily materials, energy and labour. The models also included relevant leakage factors from production, transport and storage. In generic terms, the LCOC is calculated as
\begin{equation}
\text{LCOC} = w + f + c = w + \frac{\sum{PV(F_t)}}{\sum{CDR_t}} + \frac{SP}{\sum{CDR_t}}*CRF
\end{equation}
where $w$ is the variable cost per tonne of CDR, $PV(F)$ is the present value of the annual fixed costs for operation and maintenance, $SP$ is the system price (i.e., the capital investment) required for the project, and $CRF$ is the capital recovery factor, which, at an assumed discount rate of 7\% and project lifetime of 20 years, is set at 9.44\%. For the other LCOC-derived terms, such as the LCOCR, the costs of transportation and storage are added, co-revenues deducted, and the leakage factor applied to the sum of CDR.\vspace{2mm}

For the input data, country-level data on the availability and, if relevant, cost of the primary constraints, such as land and biomass, was gathered primarily from the FAOSTAT database \parencite{fao_faostat_2025, fao_faostat_2025-1, fao_faostat_2026}. Additionally, country-level data on the cost of labor, materials and other relevant inputs was collected. If such data was not available, global averages were calculated and adjusted for each country using purchasing power parity (PPP) data of the World Bank \parencite{world_bank_icp_2021}. The collected data was then fed into the techo-economic models in order to build a global marginal abatement cost curve (MACC) for each NET.

\include*{afrf}
\include*{bicrs}
\include*{ewr}
\include*{dac}

\section{Overview of modeling results}
As a summary answer to research question one, Table 3.5 provides an overview of the CDR potential and LCOC ranges of the four approaches modeled in this chapter. Cost and CDR potential generally fall within or close to the ranges of the previous literature (see table 2.2), with the exception of ERW, which was found to be available at a minimum cost about 50\% below the literature estimates.
\begin{table}[h]
%\centering
\small
\renewcommand{\arraystretch}{1.4}
\begin{tabular}{lrrr}
Technology     & \begin{tabular}[c]{@{}l@{}}Lowest LCOC\\ (USD/t CO$_2$)\end{tabular} & \begin{tabular}[c]{@{}l@{}}Highest LCOC\\ (USD/t CO$_2$)\end{tabular} & \begin{tabular}[c]{@{}l@{}}Potential\\ (Gt CO$_2$ / yr)\end{tabular} \\ \hline
\textbf{A/R}   & 3                                                                    & 88                                                                    & 2,431                                                                \\
\textbf{BiCRS} & 116                                                                  & 319                                                                   & 4,793                                                                \\
\textbf{ERW}   &                                                                      &                                                                       & 8,630                                                                \\
\textbf{DAC}   & 328                                                                  & 527                                                                   & -                                                                    \\ \hline
\end{tabular}
\caption{Modeled cost and potential of selected NET}
\end{table}

\mychapter{4}{Development of optimal global CDR portfolios}
To construct the ideal portfolios, two types of optimization approaches were considered. First, a simple static portfolio was created, which is based on a one-off removal of a given quantity of carbon dioxide on a given timeline, such as the removal of ten gigatonnes of CO$_2$ over five years. In the second step, this simple static approach was evolved into a dynamic portfolio approach, which considers a changing (increasing) annual requirement of CDR in order to reach certain targets on a certain timeline, such as the annual removal of 10 gigatonnes of CO$_2$ by the year 2050.\vspace{2mm}

For both types, the techno-economic models described in section three were implemented in the Python programming language, using the Pyomo package and the Gurobi solver. For the portfolio optimization, multi‑period linear programs (LPs with SOS2) were used, for which the objective was to minimize the total costs, constrained by the availability of each NET and the requirement of fulfilling the respective CDR target.\vspace{2mm}

The BECCS model is used as the representative for both BiCRS approaches. This simplification is valid because both BiCRS approaches are constrained by, and therefore compete for, the same biomass feedstock base, and the marginal abatement cost curves of both technologies exhibit a similar cost and availability range. As a result, explicitly including both models would primarily affect the allocation within the BiCRS category, but not materially impact the overall portfolio results. This modeling choice reflects a complexity-conscious implementation of the BiCRS pathways, but does not imply that BECCS and biochar are identical on a technological level. The cost of ERW is modeled using a piecewise linear function derived from the MACC, as created above, in order to allow the model to choose any operating level and compute the resulting total costs without introducing nonlinearities. Since the capital expenditures of ERW make up less than 2\% of its LCOC, using the MACC as a fixed, exogenous input, instead of separately modeling the capital investments, is a sufficiently realistic approximation. The same appraoch was used for BECCS in the static models. For the dynamic BECCS model, the fixed and variable costs and the capital investments were modeled separately, meaning that the optimizer chooses not only the operating level for BECCS, but also the timeline for capacity construction. Piecewise functions were also applied for the rock use of ERW and the learning curve of DAC. In Pyomo, they were implemented via SOS2 piecewise constraints, which approximate nonlinear cumulative cost curves. The CapEx for DAC are defined by a bilinear function of the amount built in a given year and the learning-curve based cost per unit. A McCormick envelope \parencite{mccormick_computability_1976} was used as a linear approximation for this function by introducing an auxiliary variable for CapEx and four linear constraints with variable bounds. The A/R model is linear and does not require the use of piecewise approximation.\vspace{2mm}

The following restrictions were considered in the optimization models: A/R is restricted by the growth of trees, the total available land area, and a limitation on the annual reforested area relative to all available land. In addition, an annual reversal rate is implemented to simulate fires and pests. The BECCS model is limited by the available amount of waste biomass, an initial availability of electric power capacity, and an annual maximum capacity increase of, while the technology itself is assumed to be available at a constant construction cost. ERW is restricted by the agricultural land area available for spreading, as well as a limitation of annual mining capacity relative to the global coal mining industry, which increases linearly over the operating years. DAC is only limited by an assumed initial capacity and a maximum annual capacity increase of relative to the previous year's capacity.

\include*{stat_port}
\include*{dyn_port}

\section{Discussion}

\subsection{Static Portfolios}

From the static portfolios shown in figures 4.1 and 4.2, a number of conclusions can be drawn. First of all, the target formulation has a significant impact on the results. With the annual constraint, the removals must be delivered immediately, which forces the use of more costly BECCS removals to cover the gap left by cheaper but slower A/R and ERW. Because A/R has annual planting limits and a growth lag, it is not able to fully satisfy short-term targets on its own, especially not in the case of annual constraints. Enhanced weathering faces similar constraints regarding its expansion and use of rock material, while BECCS is deployable much quicker. A counter-intuitive case seems to arise in the annually-constrained model at the 20 and 50 year horizons, which use less ERW than the shorter-term scenarios while still deploying some BECCS. However, looking at figure A.5 in the appendix, one can see that this is caused by the BECCS deployment in year one, which is required to fulfill that year's target before either of the other approaches can become effective. With the overall removal constraint, the optimizer can delay removals from the first year to later times. This temporal flexibility allows it to rely nearly entirely on the cheaper A/R approach, and makes engineered CDR largely disappear for the longer time horizons. Lastly, it is worth noting that DAC sees no application in either of the scenarios due to its high costs compared to the other NET's, making it non-competitive under the static optimization model.\vspace{2mm}

The stark difference in costs between the long and short time-frames is similarily explained not only by the increased use of cheaper A/R in the longer time-frames, but also the fact that A/R itself becomes cheaper the longer the trees are allowed to grow. This is due to the fact that A/R is largely defined by its upfront costs for forest restoration, while carbon removal occurs over the entire growth cycle of 30 to 50 years. In the longer time-frames, each tree therefore removes more carbon, making A/R more efficient and decreasing its LCOC by an order of magnitude.\vspace{2mm}

From the perspective of VCM participants, this makes A/R the most attractive approach. For example, if a company wanted to achieve the goal of "historic carbon neutrality" by removing the sum of its past emissions in a one-off project, the model suggests that simply buying and retiring carbon removal credits from afforestation is the most economical solution.\vspace{2mm}

These results illustrate the limitations of a simple cost-based approach to selecting negative-emission technologies discussed in chapter 2 of this thesis. Under these static constraints, low-cost removals dominate the solution, especially if temporal flexibility is available, while novel NET are used only to fill near-term gaps. The BECCS and EWR approaches serve as temporary "compliance technologies" and Direct Air Capture is not deployed at all. This result carries two issues common for such a static, present-cost based perspective: First, the potential cost reduction due to increased deployment, and the connected learning effects especially of DAC, is entirely disregarded. However, even if cost reductions were considered, A/R would likely remain the cheapest option. The second, and more fundamental, issue is that this type of model treats the low-cost attractiveness of A/R as if it was indefinitely available, ignoring the limited nature of A/R, which is in reality constrained by the global land area available for forest restoration. Once this limited potential is exhausted, reliance on A/R alone is no longer possible. Early investment in engineered NET may therefore be necessary not because they are currently the cost-optimal option, but to drive down future costs and to ensure long-term CDR availability once cheaper approaches hit their natural limits.

\subsection{Dynamic Portfolio}

The dynamic portfolio shows staggering numbers, both, in regard to the total amount of carbon removed until the year 2100, and in regard to the total cost necessary to achieve this target. For comparison, the 855 Gt of CO$_2$ requiring removal over the next 75 years are approximately equivalent to the entire historic emission of humanity until 1990, while the sum of emissions since 1990 amounts to about 1000 Gt \parencite{ritchie_co_2020}.\vspace{2mm}

The development of the LCOC of the various technologies also provide intersting insights. The nearly flat cost of ERW is caused largely by the modeling constraints, which limit the actual deployment of the technology to a point significantly below the hypothetical potential shown by the MACC, which was developed in section 3.3. The increasing LCOC of A/R, and to a more significant degree, of BECCS, are explained by the fact that the optimizer uses the cheapest potential first, e.g. the tropic rainforests with very fast growth and high biomass accumulation, or the cheaply available manure for BECCS. Once these options are are exhausted, it moves to add the more expensive options from less economical areas or input materials. By the mid-2040s, both, A/R and BECCS, will reach their maximum annual potential. This means that while they will keep contributing significant amounts of CDR, their expansion will be stopped by natural limitations in land and waste biomass availability. Finally, the results show that A/R removals start to decline from the mid-2060s, when all available land is used up and the maturing forests start to absorb less and less CO$_2$ from the atmosphere.\vspace{2mm}

Looking at DAC, the fact that its LCOC seems to only decreases slightly despite the sharp drop in the capital expenditures per unit of capacity can be traced to the nature of the learning effects: The amount of DAC capacity constructed under the early, high-cost conditions makes up only a small share of the total DAC capacity. By the time the 100th Mt of capacity is deployed, the cost will already have fallen to below 1000 USD, and by the first Gt, it will be below 500 USD. Additionally, DAC construction begins earlier than the LCOC curve, which only starts from the point where DAC commences operations. Investing in such plants, only for them to sit idle for nearly their entire lifetime, seems counterintuitive. However, it makes sense considering that operating them is still more expensive than using the other NET as long as these can fulfill the annual CDR target. The initial DAC investment therefore occurs largely to build up capacity to be ready when it is required, while "buying down" the average capital expenditures per tonne of capacity. Therefore, these results should be regarded as an encouragement to begin DAC deployment early in order to drive down its costs by the time it is required in significant amounts from the mid-2040s onwards.\vspace{2mm}